\documentclass[12pt]{article}
\usepackage[T1]{fontenc}
\usepackage[utf8]{inputenc}
\usepackage{lmodern}
\usepackage{textcomp}
\usepackage{enumitem}
\usepackage{geometry}
\geometry{margin=1in}
\begin{document}
\begin{center}
Answer Key - Psychology - Graduate Level Assessment 17 \
\small (Answers are ordered exactly as in the question paper.)
\end{center}

\section*{Multiple Choice}
\begin{enumerate}[label=\arabic*.]
\item \textbf{Answer:} A
\\ \textit{Options:} A) 18 to 20; B) 5 to 7; C) 13 to 15; D) 10 to 12; E) 3 to 4
\\ \textit{Note:} Children begin to deliberately and regularly use rehearsal, elaboration, and organization as memory strategies around 18 to 20 years of age, as this period aligns with the development of advanced cognitive processes and metacognitive awareness that facilitate effective learning and memory retention.
\item \textbf{Answer:} B
\\ \textit{Options:} A) conversion disorder.; B) schizophrenia.; C) OCD.; D) post-traumatic stress disorder.; E) bipolar disorder.
\\ \textit{Note:} Delusions of grandeur, which involve an inflated sense of self-importance and power, are a common symptom of schizophrenia, reflecting the disorder's impact on perception and reality.
\item \textbf{Answer:} D
\\ \textit{Options:} A) the understanding of cause and effect.; B) the capacity for deductive reasoning.; C) the ability to reason about abstract concepts and problems.; D) the symbolic function.; E) the object concept.
\\ \textit{Note:} The correct answer is the symbolic function because during Piaget's preoperational stage, children develop the ability to represent objects and experiences through symbols, which is fundamental for cognitive development and imaginative play.
\item \textbf{Answer:} A
\\ \textit{Options:} A) One of the family members has severe depression.; B) Family members are dealing with the loss of a loved one.; C) Family members attribute their problems to one of the members.; D) The presenting problem involves sibling rivalry.; E) The family has recently moved to a new city.
\\ \textit{Note:} Family therapy may be contraindicated when one family member has severe depression because their individual mental health issues can disrupt the therapeutic process and may require focused, individualized treatment rather than a family-oriented approach.
\item \textbf{Answer:} A
\\ \textit{Options:} A) Consciousness is the storage of genetic information in cells.; B) Consciousness is the ability to think; C) Consciousness is the ability to talk to oneself; D) Consciousness is the exclusive ability to solve complex mathematical problems.; E) Consciousness is the biological process of breathing and heartbeat regulation.
\\ \textit{Note:} The correct answer is inaccurate because consciousness, in the context of psychology, refers to the state of being aware of and able to think and perceive one's surroundings, rather than the storage of genetic information in cells, which pertains to biological functions.
\end{enumerate}

\section*{True / False}
\begin{enumerate}[label=\arabic*.]
\setcounter{enumi}{5}
\item \textbf{Answer:} False
\\ \textit{Options:} A) True; B) False
\\ \textit{Note:} The statement is false because the resting potentials of retinal and cochlear receptors differ due to their distinct cellular environments and the specific ionic compositions of their respective endolymph and perilymph fluids.
\item \textbf{Answer:} True
\\ \textit{Options:} A) True; B) False
\\ \textit{Note:} In Milgram's obedience experiments, a significant majority of participants, approximately 60\%, were convinced that they had administered the maximum shock level to the learner, illustrating the powerful influence of authority on individual behavior and the perception of their own actions.
\item \textbf{Answer:} True
\\ \textit{Options:} A) True; B) False
\\ \textit{Note:} Abraham Maslow's hierarchy of needs posits that individuals must satisfy lower-level needs, such as esteem needs, including self-esteem and recognition, before they can pursue higher-level needs related to self-actualization, which involves realizing one's full potential and personal growth.
\item \textbf{Answer:} False
\\ \textit{Options:} A) True; B) False
\\ \textit{Note:} Personality and intelligence are distinct constructs, with research indicating that while they may interact, cognitive abilities are influenced by a variety of factors including genetics, environment, and education, rather than being solely determined by personality traits.
\item \textbf{Answer:} True
\\ \textit{Options:} A) True; B) False
\\ \textit{Note:} The statement is true because psychology relies on empirical evidence gathered through systematic observation and measurement techniques to analyze and understand behavior and mental processes.
\end{enumerate}

\section*{Long Form Response}
\begin{enumerate}[label=\arabic*.]
\setcounter{enumi}{10}
\item \textbf{Answer:} Basic somatosensations, such as touch, pain, temperature, and proprioception, play a critical role in human perception by providing essential information about the body's interaction with the environment. These sensations are fundamental for survival and facilitate behavioral responses to external stimuli. In contrast, itch is not classified among these basic somatosensations due to its distinct neurophysiological pathways and subjective experiences; it is often viewed as a complex sensation that can reflect both physiological and psychological states. This distinction has significant implications for psychological research, as it necessitates a differentiated approach to studying sensory experiences and addressing conditions like chronic itch, which may involve both somatic and emotional components. Understanding itch's unique characteristics can enhance our grasp of sensory processing and the interplay between physical sensations and psychological well-being.
\\ \textit{Note:} correct because it accurately differentiates basic somatosensations, which are critical for survival and immediate behavioral responses, from itch, which has unique neurophysiological mechanisms and is considered a more complex sensory experience, highlighting its implications for psychological research on sensory perception.
\item \textbf{Answer:} To analyze the home run data of the Boston Red Sox over eight consecutive games, we calculate the mean, median, and mode, all of which are found to be 3. The mean, representing the average number of home runs per game, indicates a consistent level of performance across the games. The median, which is the middle value when the data is ordered, reinforces the idea of stability in their performance, suggesting that half of the games resulted in home runs equal to or below this number. Finally, the mode, the most frequently occurring value, also being 3, reflects that this performance level is not only typical but perhaps a target for the team. Together, these statistical measures provide valuable insights into the team's performance trends, highlighting reliability in their offensive output and potentially guiding strategies for future games.
\\ \textit{Note:} correct because it demonstrates that the mean, median, and mode being equal to 3 highlights a consistent and stable performance in home runs by the Boston Red Sox over the eight games, reflecting a reliable trend in their offensive capabilities.
\end{enumerate}

\vfill
\noindent\textit{End of Answer Key}
\end{document}